% ========================================================================
%  DROP-IN REPLACEMENT for paper.tex lines 99-170
%  (\section{Introduction} through the end of \section{Preliminaries})
%
%  Preamble additions required (place next to \newcommand{\Hom}):
%
%    \DeclareMathOperator{\Sp}{Sp}          % spider tree
%    \newcommand{\Aut}{\mathrm{Aut}}
%    \newcommand{\Wt}[1]{w_{#1}}            % walk vector
%    \newcommand{\ip}[2]{\langle #1, #2 \rangle_{\mathbf{s}}}
%    \usepackage{amsthm}                    % already loaded
%
%  No new packages are needed.
% ========================================================================

% ════════════════════════════════════════════════════════════════════════
\section{Introduction}
\label{sec:intro}

\subsection{Counting maps between graphs}

Let $G$ and $H$ be finite graphs, where $H$ may carry loops. A
\emph{homomorphism} $\varphi \colon G \to H$ is a map
$V(G) \to V(H)$ that sends edges to edges: if $uv \in E(G)$ then
$\varphi(u)\varphi(v) \in E(H)$. We write $\Hom(G,H)$ for the number of
such maps.

The quantity $\Hom(G,H)$ is a single number that specialises to many
familiar ones. If $H = K_q$ is the complete graph, $\Hom(G,K_q)$ counts
proper $q$-colourings of $G$; if $H$ is a single edge with a loop at one
end, $\Hom(G,H)$ counts independent sets; if $H$ is a path, $\Hom(G,H)$
counts height functions with bounded increments. For this reason $H$ is
called the \emph{target} and the map count is often called the number of
\emph{$H$-colourings} of $G$.

Two extremal questions present themselves at once. Among all trees $T$
on $n$ vertices, which maximises $\Hom(T,H)$, and which minimises it?
The first question is settled: Sidorenko~\cite{sidorenko1991} showed
that the star $K_{1,n-1}$ is always a maximiser, for every target $H$.
The second is not.

\subsection{The path conjecture and its failure}

Galvin and Tetali~\cite{galvin-tetali} proposed the natural companion:
that the path $P_n$ should always be a minimiser. The intuition is
appealing. A homomorphism from a tree is built by propagating
constraints outward from a root, and the path is the tree that spreads
those constraints along the longest possible chain, giving each vertex
the fewest independent choices. Csikv\'ari and Lin~\cite{csikvari-lin}
developed the systematic tool for this circle of ideas, the
KC-transformation, and proved the conjecture for large classes of
targets.

The conjecture is nevertheless false. Leontovich~\cite{leontovich1989}
exhibited a target for which some tree beats the path, and
Galvin, McMillon, Redlich and Nir~\cite{galvin2025} gave explicit
bipartite counterexamples and named the phenomenon. Following them we
call $H$ a \emph{Leontovich graph} when the path fails to minimise, and
we call the smallest $n$ at which it fails the \emph{crossover
threshold}. A companion paper~\cite{galvin2026} introduces the dual
class of \emph{Hoffman--London graphs}, the targets for which the path
does minimise, and poses the quantitative questions this paper takes up:
how small can a Leontovich graph be, and how early can a crossover
happen?

What makes these questions hard is that the failure is faint. In every
known example the two counts agree to within a relative margin of
$10^{-4}$ or less, and the winning tree differs from the path by a
single misplaced pendant vertex. The effect is a near-cancellation
between two exponentially growing quantities, so it is invisible to
crude estimates and easily destroyed by floating-point noise. Much of
the work below consists of building comparisons that survive at that
precision, and of saying plainly which of our conclusions are proofs
and which are the results of bounded searches.

\subsection{What this paper does}

We attack the problem from three directions.

\emph{Analytically.} We isolate a condition on the eigenvalues of a
target that forbids a crossover outright
(Theorem~\ref{thm:spectral}), which disposes of every star and every
$3$-orbit symmetric tree in one stroke. We determine the exact
crossover threshold of the smallest previously known Leontovich tree,
and show that once the sign flips it never flips back
(Theorem~\ref{thm:threshold}). We prove that passing to the bipartite
double cover preserves both the Leontovich property and its asymptotic
strengthening (Theorem~\ref{thm:double_cover}), which converts looped
witnesses into simple ones.

\emph{Computationally.} We sweep large families of targets --- all
connected simple graphs on at most $11$ vertices, all connected
bipartite graphs on at most $15$ vertices, all symmetric trees below
$78$ vertices --- against the near-path source trees defined below, and
we record the results as bounded sweeps rather than as theorems.

\emph{By synthesis.} Where enumeration runs out, we encode the
crossover condition as an integer constraint problem and hand it to an
SMT solver. This produces the two smallest bipartite witnesses we know
of, on $15$ and $18$ vertices.

\subsection{Standards of evidence}
\label{sec:evidence}

Because the arguments below mix proof with computation, and because
several claims in earlier drafts of this work did not survive contact
with exact arithmetic, we label every substantive assertion with one of
four tags and use them consistently.

\begin{description}[leftmargin=2.4em,style=nextline,font=\normalfont\bfseries]
  \item[(P) Proved.] An unconditional proof, quantifiers
    included. Stated as \textbf{Theorem}, \textbf{Proposition},
    \textbf{Lemma} or \textbf{Corollary}.
  \item[(E) Exact finite computation.] A specific numerical claim about
    named objects, verified in exact integer or certified
    arbitrary-precision arithmetic. A witness of this kind is as
    reliable as the code that produced it; the scripts are named in
    place and the values are reproducible.
  \item[(S) Bounded sweep.] A negative result of the form ``no member of
    family $\mathcal{F}$ exhibits behaviour $B$ within the tested
    window''. Such a statement is a certificate about the window, not
    about the family. Stated as \textbf{Proposition (computational
    sweep certificate)} and never as a nonexistence theorem.
  \item[(N) Certified numerics.] A high-precision floating-point
    computation accompanied by an explicit error bound --- here a
    Bauer--Fike residual bound on the Perron eigenvalue and a
    Davis--Kahan bound on its eigenvector --- so that the reported
    inequality holds with a proven margin.
\end{description}

The distinction between (P) and (S) is the one that matters most. A
filter that tests the path against near-path trees only can certify that
\emph{those} trees fail to beat the path; it cannot certify that no tree
does. We therefore never upgrade a sweep to a nonexistence theorem, and
where an earlier claim of ours did so we say where the boundary now
lies.

\subsection{Summary of contributions}

\begin{enumerate}[label=(\roman*)]
  \item \textbf{(E)} Verification of the path-minimiser theorem for the
    odd paths $H \in \{P_3,P_5,P_7,P_9\}$ over all source trees on
    $n \le 22$ vertices, and an explanation of the accompanying
    transition: for $P_3$ and $P_5$ the path is a minimiser but not a
    unique one, and the tie classes are exponentially large; from $P_7$
    onward the minimiser is unique. The mechanism is an arithmetic one
    --- $\lambda_1^2$ is rational exactly for $k \in \{1,2,3,5\}$ ---
    and it answers Problem~6.3 of~\cite{galvin2026}
    (Section~\ref{sec:odd-path}).
  \item \textbf{(P)} An eigenvalue criterion (Theorem~\ref{thm:spectral}):
    a rooted bipartite target whose active quotient has a single
    positive eigenvalue cannot be Leontovich. This rules out all
    $2$-orbit and $3$-orbit symmetric trees without any search
    (Corollary~\ref{cor:2_and_3_orbit_trees}).
  \item \textbf{(P)} The crossover threshold of $T(7,1,9)$ is exactly
    $n = 13$, and it is permanent: the sign of the difference changes
    once and never again (Theorem~\ref{thm:threshold}). The companion
    impossibility at $n = 5$ for the whole family $T(x,y,z)$ is a
    polynomial identity (Proposition~\ref{prop:n5}).
  \item \textbf{(S)} All $1{,}018{,}680{,}329$ connected simple graphs on
    $m \le 11$ vertices, and all $608{,}930{,}559$ connected bipartite
    graphs on $m \le 15$ vertices, tested against the near-path family
    with $n \le 200$, $d \le 20$. No simple target below $12$ vertices
    produced a near-path crossover. This bounds, but does not settle,
    Problem~4.3 of~\cite{galvin2025} (Section~\ref{sec:bounds}).
  \item \textbf{(E)} Two small bipartite witnesses: a $15$-vertex graph
    $H^*$ whose earliest crossover is at $(n,d) = (49,16)$, and an
    $18$-vertex graph $H_{18}$ that crosses over against the classical
    depth-$2$ near-path at $n = 17$ with margin
    $\Delta = 5{,}068{,}778$ (Section~\ref{sec:discussion}).
  \item \textbf{(P)} The bipartite double cover preserves both the
    Leontovich and the strongly Leontovich property
    (Theorem~\ref{thm:double_cover}). Applied to the looped tree
    $\hat T(1,35,1,50)$, this yields a \emph{simple} strongly Leontovich
    bipartite graph on $3{,}644$ vertices, whose first even crossover we
    locate exactly at $n = 17{,}340$ --- evidence that even-$n$
    crossovers do not require a loop, only enough room to delay the
    cancellation.
  \item \textbf{(N)} A perturbed double cover on $6{,}806$ vertices whose
    certified leading-coefficient ratio is $0.999952140676 < 1$,
    a candidate simple non-bipartite strongly Leontovich target.
\end{enumerate}

Section~\ref{sec:prelim} fixes notation and records the two
transfer-matrix identities that every computation in the paper rests on.

% ════════════════════════════════════════════════════════════════════════
\section{Preliminaries}
\label{sec:prelim}

Throughout, $H$ denotes a finite connected target graph, possibly with
loops, on $m = |V(H)|$ vertices, with adjacency matrix $A = A(H)$; a
loop at $u$ contributes $A_{uu} = 1$. Source graphs are trees on $n$
vertices. All logarithms are natural and all matrices are real.

\subsection{Walks are the only thing being counted}

The single fact underlying every formula below is elementary. Let
$\mathbf{1}$ be the all-ones vector and set
\[
  w_k \;=\; A^k \mathbf{1},
  \qquad\text{so that}\qquad
  w_k(u) \;=\; \#\{\text{walks of length } k \text{ in } H
  \text{ starting at } u\}.
\]
A homomorphism from a tree is determined by choosing an image for a root
and then extending independently along each branch, so every count we
need is a sum of products of entries of the $w_k$.

\begin{lemma}[Path and spider counts]
  \label{lem:transfer}
  Let $\Sp(a_1,\dots,a_r)$ denote the \emph{spider}: the tree obtained by
  identifying one endpoint of $r$ paths with $a_1, \dots, a_r$ edges
  respectively at a common centre. Then
  \[
    \Hom(P_n, H) \;=\; \mathbf{1}^{T} A^{\,n-1} \mathbf{1}
      \;=\; \sum_{u \in V(H)} w_{n-1}(u),
    \qquad
    \Hom\bigl(\Sp(a_1,\dots,a_r), H\bigr)
      \;=\; \sum_{u \in V(H)} \prod_{i=1}^{r} w_{a_i}(u).
  \]
\end{lemma}

\begin{proof}
  Root the spider at its centre and condition on the image $u$ of the
  root. The branches are vertex-disjoint apart from the centre, so their
  images may be chosen independently; the choices along a branch with
  $a_i$ edges are exactly the walks of length $a_i$ from $u$. The path
  is the case $r = 2$, or directly the $k$-step iteration of $A$.
\end{proof}

\subsection{The near-path family}

\begin{definition}[Near-path trees]
  \label{def:nearpath}
  For $n \ge 5$ and $1 \le d \le n-3$, the \emph{near-path tree} is the
  spider
  \[
    E_n^{(d)} \;=\; \Sp\bigl(n-d-2,\; d,\; 1\bigr),
  \]
  the tree on $n$ vertices obtained from the path $P_{n-1}$ by attaching
  a pendant vertex to the vertex at distance $d$ from one end. We write
  $E_n = E_n^{(2)}$ for the classical case, the tree used by
  Leontovich~\cite{leontovich1989}, obtained by hanging a pendant on the
  third-from-last vertex of $P_{n-1}$.
\end{definition}

Combining Definition~\ref{def:nearpath} with Lemma~\ref{lem:transfer},
\begin{equation}
  \label{eq:nearpath-count}
  \Hom\bigl(E_n^{(d)}, H\bigr)
    \;=\; \sum_{u \in V(H)} w_{n-d-2}(u)\, w_{d}(u)\, w_{1}(u),
\end{equation}
which is the identity implemented by every verification script in this
work. Two consequences are worth recording.

\begin{remark}[The depth parameter is a leg length]
  \label{rem:d-symmetry}
  A spider does not remember the order of its legs, so
  $E_n^{(d)} \cong E_n^{(n-2-d)}$. The parameter $d$ is therefore not a
  depth so much as the length of the shorter of the two long legs, and a
  sweep over $d \le d_{\max}$ covers exactly those near-path trees with
  $\min\{d,\, n-2-d\} \le d_{\max}$. Our sweeps use
  $d_{\max} = 20$ with $n \le 200$; the untested trees are those whose
  pendant sits near the middle of the path.
\end{remark}

\begin{remark}[Why this family and not another]
  \label{rem:why-nearpath}
  The near-path trees are the natural first place to look, for a reason
  made precise in Remark~\ref{prop:pf}: any tree that beats the path must
  match its exponential growth rate $\lambda_1^{\,n}$ and win on the
  constant in front. Branching in the interior raises the growth rate
  outright, so a challenger must confine its defect to the boundary.
  This is a heuristic for choosing a search family, not a theorem, and
  we treat every negative result obtained with it accordingly.
\end{remark}

\subsection{Symmetric trees and their quotients}

\begin{definition}[Spherically symmetric trees]
  \label{def:symtree}
  For $d_1,\dots,d_k \ge 1$, the tree $T(d_1,\dots,d_k)$ is rooted, of
  depth $k$, in which every vertex at depth $i-1$ has exactly $d_i$
  children. Its vertex set falls into $k+1$ \emph{orbits} by depth, of
  sizes $s_0 = 1$ and $s_i = d_1 d_2 \cdots d_i$, and
  \[
    |V| \;=\; 1 + d_1 + d_1 d_2 + \cdots + d_1 d_2 \cdots d_k .
  \]
  We write $\hat T(d_1,\dots,d_k)$ for the same tree with a loop added at
  the root. The case $k = 3$ recovers the classical family
  $T(x,y,z)$ on $1 + x + xy + xyz$ vertices, whose four orbits are root,
  children, grandchildren and leaves.
\end{definition}

The loop matters: it destroys bipartiteness, and with it the parity
obstruction of Section~\ref{sec:leontovich}.

\begin{definition}[Equitable partition and quotient matrix]
  \label{def:quotient}
  A partition $\pi = \{C_1,\dots,C_K\}$ of $V(H)$ is \emph{equitable} if
  there are integers $q_{ij}$ such that every vertex of $C_i$ has exactly
  $q_{ij}$ neighbours in $C_j$. The matrix $Q = (q_{ij})$ is the
  \emph{quotient matrix} of $\pi$ and $\mathbf{s} = (|C_1|,\dots,|C_K|)$
  its \emph{size vector}. The orbit partition of any automorphism group
  is equitable; so is the coarsest partition produced by colour
  refinement.
\end{definition}

\begin{lemma}[Reduction to the quotient]
  \label{lem:quotient}
  Let $\pi$ be equitable with quotient $Q$ and size vector $\mathbf{s}$.
  Then each walk vector $w_k$ is constant on cells, and its cell values
  are the entries of $Q^k \mathbf{1}_K$. Consequently, for every $n$ and
  every spider,
  \[
    \Hom(P_n,H) \;=\; \mathbf{s}^{T} Q^{\,n-1}\mathbf{1}_K,
    \qquad
    \Hom\bigl(\Sp(a_1,\dots,a_r),H\bigr)
      \;=\; \sum_{i=1}^{K} s_i \prod_{j=1}^{r} \bigl(Q^{a_j}\mathbf{1}_K\bigr)_i .
  \]
\end{lemma}

\begin{proof}
  Induct on $k$. The claim holds for $w_0 = \mathbf{1}$. If $w_{k}$ takes
  the value $x_i$ on $C_i$, then for $u \in C_i$,
  $w_{k+1}(u) = \sum_j q_{ij} x_j$, which depends only on $i$. Summing
  Lemma~\ref{lem:transfer} cell by cell gives the two displayed
  identities.
\end{proof}

Lemma~\ref{lem:quotient} is what makes the computations in this paper
feasible: a target on $6{,}806$ vertices is replaced by an $18 \times 18$
integer matrix with no loss of information about $\Hom(T,\cdot)$ for the
trees we test. It also explains the recurring phrase \emph{active
spectrum}: only the eigenvalues of $Q$, not those of the full adjacency
matrix, can appear in these counts. The eigenvalues of $A(H)$ orthogonal
to the quotient subspace are invisible here, and we say they are
\emph{inert}.

\begin{lemma}[The quotient is self-adjoint in the size weighting]
  \label{lem:selfadjoint}
  Define $\ip{x}{y} = \sum_i s_i x_i y_i$. Then $Q$ is self-adjoint with
  respect to $\langle\cdot,\cdot\rangle_{\mathbf{s}}$; in particular its
  eigenvalues are real and it has a $\langle\cdot,\cdot\rangle_{\mathbf{s}}$-orthonormal
  eigenbasis.
\end{lemma}

\begin{proof}
  Counting edges between $C_i$ and $C_j$ in two ways gives
  $s_i q_{ij} = s_j q_{ji}$, i.e.\ $\ip{Qx}{y} = \ip{x}{Qy}$.
  Equivalently $S^{1/2} Q S^{-1/2}$ is a symmetric matrix, where
  $S = \operatorname{diag}(\mathbf{s})$.
\end{proof}

\subsection{Leading coefficients and the two Leontovich properties}

Let $\lambda_1 \ge \lambda_2 \ge \cdots \ge \lambda_K$ be the eigenvalues
of $Q$ with $\langle\cdot,\cdot\rangle_{\mathbf{s}}$-orthonormal
eigenvectors $u^{(1)},\dots,u^{(K)}$, and put
$\alpha_j = \ip{\mathbf{1}}{u^{(j)}}$. Lemma~\ref{lem:quotient} and
Lemma~\ref{lem:selfadjoint} give the exact expansions
\begin{equation}
  \label{eq:expansion}
  \Hom(P_n,H) = \sum_{j=1}^{K} \alpha_j^{2}\, \lambda_j^{\,n-1},
  \qquad
  \Hom\bigl(E_n^{(d)},H\bigr)
    = \sum_{j=1}^{K} \alpha_j \beta_j^{(d)}\, \lambda_j^{\,n-d-2},
  \qquad
  \beta_j^{(d)} = \ip{w_1 \odot w_d}{u^{(j)}},
\end{equation}
where $\odot$ is the entrywise product. Everything in this paper is a
statement about the interference between the $j = 1$ term and the rest.

\begin{definition}[Leontovich, strongly Leontovich]
  \label{def:leontovich}
  A target $H$ is \emph{Leontovich} if there is an $n$ and a tree $T$ on
  $n$ vertices with $\Hom(T,H) < \Hom(P_n,H)$: the path is not a
  minimiser. It is \emph{strongly Leontovich}~\cite{galvin2026} if
  moreover a fixed near-path family beats the path for all large $n$,
  i.e.\ $\Hom(E_n^{(d)},H) < \Hom(P_n,H)$ for all sufficiently large $n$
  of the relevant parity.
\end{definition}

If $H$ is connected and non-bipartite, $\lambda_1$ strictly dominates
$|\lambda_j|$ for $j \ge 2$, and \eqref{eq:expansion} identifies the
strong property with a single inequality between leading coefficients.

\begin{definition}[Leading-coefficient ratio]
  \label{def:rho}
  For a connected non-bipartite target $H$ with active quotient
  $(Q,\mathbf{s})$, Perron eigenvalue $\lambda_1$ and positive Perron
  eigenvector $u = u^{(1)}$, set
  \begin{equation}
    \label{eq:rho}
    \rho_d(H) \;=\;
      \lim_{n \to \infty}
      \frac{\Hom\bigl(E_n^{(d)},H\bigr)}{\Hom(P_n,H)}
    \;=\;
      \frac{\ip{w_1 \odot w_d}{u}}
           {\lambda_1^{\,d+1}\, \ip{\mathbf{1}}{u}} .
  \end{equation}
  The right-hand side is independent of the normalisation of $u$. Then
  $H$ is strongly Leontovich for depth $d$ precisely when
  $\rho_d(H) < 1$, and $\rho_d(H) > 1$ forbids it.
\end{definition}

\begin{remark}[Bipartite targets and parity]
  \label{rem:bipartite-parity}
  If $H$ is bipartite then $-\lambda_1$ is also an eigenvalue of $Q$ and
  the limit in \eqref{eq:rho} need not exist: the two dominant modes
  reinforce at one parity of $n$ and cancel at the other. This is not a
  technicality but the governing phenomenon. It is why every crossover
  in this paper occurs at odd $n$ unless the target is very large
  (Section~\ref{sec:doublecover}), and why the looped root in
  $\hat T(\cdot)$ is the standard device for removing the obstruction.
  For bipartite $H$ we therefore read \eqref{eq:rho} along a fixed
  parity class and say so explicitly.
\end{remark}

\begin{remark}[The margins involved]
  \label{rem:margins}
  To fix expectations: the smallest witnesses in this paper win by a
  relative margin of about $3.4 \times 10^{-4}$ ($H_{18}$ at $n = 17$)
  down to $4.8 \times 10^{-5}$ (the perturbed cover of
  Section~\ref{sec:doublecover}). A double-precision evaluation of
  \eqref{eq:nearpath-count} accumulates relative error on the order of
  $n\,m\,\varepsilon \approx 10^{-12}$, which is comfortably below these
  margins but not below every margin that might exist. Floating point is
  therefore used only to \emph{flag} candidates; every reported witness
  is re-derived in exact integer arithmetic through
  Lemma~\ref{lem:quotient}, and every asymptotic claim is accompanied by
  the certified bound described in Section~\ref{sec:evidence}.
\end{remark}

\subsection{Notation summary}

\begin{table}[H]
\centering
\caption{Notation used throughout.}
\label{tab:notation}
\begin{tabular}{@{}ll@{}}
  \toprule
  Symbol & Meaning \\
  \midrule
  $\Hom(G,H)$ & number of homomorphisms $G \to H$ \\
  $m$, $n$ & order of the target $H$; order of the source tree \\
  $A$, $w_k = A^k\mathbf{1}$ & adjacency matrix of $H$; walk-count vector \\
  $P_n$ & path on $n$ vertices ($n-1$ edges) \\
  $\Sp(a_1,\dots,a_r)$ & spider with legs of $a_1,\dots,a_r$ edges \\
  $E_n^{(d)} = \Sp(n{-}d{-}2, d, 1)$ & near-path tree; $E_n = E_n^{(2)}$ \\
  $T(d_1,\dots,d_k)$, $\hat T(\cdot)$ & symmetric tree; same with a looped root \\
  $Q$, $\mathbf{s}$, $K$ & quotient matrix, cell sizes, number of cells \\
  $\lambda_1 > |\lambda_j|$, $u$ & Perron eigenvalue and eigenvector of $Q$ \\
  $\lambda_2^{*}$ & largest active subdominant eigenvalue \\
  $\rho_d(H)$ & leading-coefficient ratio, \eqref{eq:rho} \\
  $\Delta_n = \Hom(P_n,H) - \Hom(E_n^{(d)},H)$ & crossover margin; $\Delta_n > 0$ means the path loses \\
  \bottomrule
\end{tabular}
\end{table}
